\documentclass{article}

%maths
\usepackage{mathtools}
\usepackage{amsmath}
\usepackage{amssymb}
\usepackage{amsfonts}

%tikzpicture
\usepackage{tikz}
\usepackage{scalerel}
\usepackage{pict2e}
\usepackage{tkz-euclide}
\usetikzlibrary{calc}
\usetikzlibrary{patterns,arrows.meta}
\usetikzlibrary{shadows}

%pgfplots
\usepackage{pgfplots}
\pgfplotsset{compat=newest}
\usepgfplotslibrary{statistics}
\usepgfplotslibrary{fillbetween}

%colours
\usepackage{xcolor}

\usepackage{nicefrac}

\begin{document}
\section{Conic Sections}
When a double cone is sliced by a plane, the cross section is formed by the plane and the cone is called a \textbf{Conic Section}. The four main conic sections are the circle, the parabola, the ellipse, and the hyperbola.

\subsection{Circle}
A \textbf{circle} is a set of point which are equidistant from one point. The one point being called the \textbf{center} and the distance of the line is the \textbf{radius}. The standard form equation with the center at (0,0) is: $x^2 + y^2 = r^2$

\subsection{Parabola}
A \textbf{parabola} is the set of points in a plane that are the same distance from a given point and a given line in that plane.\\
\\The given point is the \textbf{focus} and the line is called the \textbf{directrix}.\\
\\The midpoint of the perpendicular segment from the focus to the directrix is called the \textbf{vertex} of the parabola.\\
\\The line which passes through the vertex and focus is called the \textbf{axis of symmetry}.
\newpage
\subsection{Ellipse}
An \textbf{Ellipse} is the set of points in a plane such that the sum of the distances from two given points in that plane stays constant.\\
\\The two points are each called a \textbf{focus point} -- plural \textit{foci}.\\
\\The midpoint of the segment joining the foci is called the \textbf{center} of the ellipse.\\
\\Ellips have two axes of symmetry: \\Longer one = \textbf{Major Axis}, Shorter one = \textbf{Minor Axis}.\\
The \textbf{major intercepts} are where the 2 points intercept with the major axis.\\
The \textbf{minor intercepts} are where the 2 points intercept with the minor axis.\\
\\

\subsection{Hyperbola}
A \textbf{hyperbola} is the set of all points in a plane such that the absolute value of the difference of the distances between two given points stays constant.\\
\\Two given points are the \textbf{foci} of the hyperbola.\\
\\Hyperbolas looks like two opposing 'U-shaped' curves.\\
\\Hyperbolas have two axes of symmetry:\\ The one going through the hyperbola's foci is the \textbf{transverse axis.}\\ The one going through the center and is perpendicular to the transverse is the \textbf{conjugate axis}.

\end{document}
